\color{unidarkgreen}{\subsection{Максимальная токовая защита стороны НН}}
\color{black}

\begin{enumerate}[label=\arabic{section}.\arabic{subsection}.\arabic*,  align=left, labelsep=4pt, leftmargin=0pt, itemindent=57pt]

\item
Максимальная токовая защита (МТЗ) стороны НН предназначена для резервирования действия
основных быстродействующих защит. Защита выполнена двухступенчатой.
Измерительные цепи защиты включены на токи фаз А, В, С со стороны НН автотрансформатора. 

\item
Каждая из ступеней имеет регулируемые значения по току и времени срабатывания. Значения
срабатывания по току определяются параметрами \textbf{«Iср МТЗ-1 НН»} и \textbf{«Iср МТЗ-2 НН»} для первой
и второй ступеней соответственно. Задержки на срабатывание защиты задаются параметрами \textbf{«Тср МТЗ-
1 НН1(НН)»} и \textbf{«Тср МТЗ-2 НН1(НН)»}.

\item
Каждая из ступеней защиты может выполнена с пуском по напряжению. Ввод контроля пуска по
напряжению в логику работы ступеней осуществляется программными накладками \textbf{«Режим КПОН МТЗ-
1 НН»} и \textbf{«Режим КПОН МТЗ-2 НН»} для первой и второй ступеней соответственно.

\item
Каждая из ступеней защиты может выполнена с блокировкой по второй гармонике (отношению
уровня второй гармоники к первой). Данная блокировка реализована на базе модуля «ОВ БНТ».
Ввод контроля броска намагничивающего тока выполняется программными накладками \textbf{«Режим БНТ МТЗ-1
НН)»}, \textbf{«Режим БНТ МТЗ-2 НН»}. 

\begin{figure}[H]
\centering
\includegraphics[width=0.75\textwidth,height=0.75\textheight,keepaspectratio]{\fbpath/053.0331 МТЗ НН/img/MTZ_NN.pdf}
\captionof{figure}{Функциональный блок МТЗ НН}\label{pic:MTZ_NN}
\end{figure}

\small
\begin{longtable}{|>{\centering\arraybackslash}m{5.3cm}|>{\centering\arraybackslash}m{4.2cm}|>{\centering\arraybackslash}m{1cm}|>{\centering\arraybackslash}m{3.3cm}|>{\centering\arraybackslash}m{1.8cm}|}
% Первая страница: полный заголовок
\caption{Перечень уставочных параметров МТЗ\hfill\vspace{-0.5\baselineskip}}\label{desc:tbl1}\\
\hline
\rowcolor{gray!30}
Наименование уставки & Диапазон & Шаг & Значение по умолчанию & Единица измерения \\
\hline
\endfirsthead

% Промежуточные страницы: «Продолжение таблицы»
\caption*{\hspace{3pt}\emph{Продолжение таблицы \ref{desc:tbl1}}\hfill\ }\vspace{-0.5\baselineskip} \\
\hline
\rowcolor{gray!30}
Наименование уставки & Диапазон & Шаг & Значение по умолчанию & Единица измерения \\
\hline
\endhead

%===m>ДЗАТ|-
%===m
\end{longtable}
\normalsize	

\end{enumerate}